% Generated draft: requires mathematical and source review.
\documentclass[11pt]{article}
\usepackage[margin=25mm]{geometry}
\usepackage{fontspec}
\setmainfont{Latin Modern Roman}
\usepackage{amsmath,amssymb,mathtools,unicode-math}
\setmathfont{Latin Modern Math}
\usepackage{enumitem,needspace,xurl,hyperref}
\setlength{\parindent}{0pt}
\setlength{\parskip}{5pt}
\setlength{\emergencystretch}{3em}
\newcommand{\N}{\mathbb{N}}
\newcommand{\Z}{\mathbb{Z}}
\newcommand{\Q}{\mathbb{Q}}
\newcommand{\R}{\mathbb{R}}
\newcommand{\C}{\mathbb{C}}
\newcommand{\F}{\mathbb{F}}
\newcommand{\A}{\mathbb{A}}
\newcommand{\PP}{\mathbb{P}}
\newcommand{\E}{\mathbb{E}}
\newcommand{\eps}{\varepsilon}
\newcommand{\dd}{\,\mathrm{d}}
\begin{document}
\section*{UnsolvedMath ID: 30001599}
\textbf{Source problem code:} catalogue entry
\subsection{Definitions and mathematical statement}
For a conditional formulation, let \(k\) be a field, let \(n,r\geq 1\) be integers, and let \(X\subseteq\PP^n_k\) be a nonempty finite set of distinct \(k\)-rational points. These parameter and rationality conventions make the statement precise, but their intended scope is not specified in the supplied source. Put \(S=k[x_0,\ldots,x_n]\). For \(P\in X\), let \(\mathfrak p_P\subset S\) be its homogeneous vanishing ideal, and define
\[
I_X=\bigcap_{P\in X}\mathfrak p_P,\qquad m=rn-n+1=n(r-1)+1.
\]
The uniform fat-point scheme of multiplicity \(m\) supported on \(X\) is
\[
Y=\operatorname{Proj}(S/I_Y),\qquad I_Y=\bigcap_{P\in X}\mathfrak p_P^{\,m}.
\]
For either scheme \(Z=X\) or \(Z=Y\), define its initial degree by
\[
\alpha(Z)=\min\{d\in\Z_{\geq0}:(I_Z)_d\neq\{0\}\}.
\]
Thus \(\alpha(Y)\) is the least degree of a nonzero homogeneous polynomial vanishing to order at least \(m\) at every point of \(X\).

The supplied conjecture asks whether, for every such configuration in the intended parameter and ground-field range,
\[
\alpha(Y)\geq r\alpha(X)+(r-1)(n-1).
\]
The displayed formulation is conditional on the conventions above. The supplied record does not determine the intended ground field or characteristic, whether geometric points or more general closed points are allowed, the precise ranges of \(n\) and \(r\), or the treatment of the empty set. These qualifications require source review before this can be treated as a definitive formulation.
\subsection{Short English statement}
The conjecture asks whether requiring multiplicity \(n(r-1)+1\) at every point forces the least degree of a vanishing hypersurface to be at least \(r\alpha(X)+(r-1)(n-1)\). The ground field, parameter ranges, and empty-set convention remain unresolved in the supplied record.
\subsection{Sources}
\begin{itemize}
\item Local UnsolvedMath registry, supplied complete source record: permanent ID \(30001599\), exact source problem code \(\text{catalogue entry}\), title \emph{Alpha Bounds for Uniform Fat-Point Schemes}. No individual proposer is supplied; the proposed year is recorded as 2010.
\item Supplied source citation: \emph{Mini-Workshop: Linear Series on Algebraic Varieties (2011)}, Oberwolfach Reports, pp. 2613--2650, DOI 10.4171/owr/2010/45. \url{https://doi.org/10.4171/owr/2010/45}.
\item Additional supplied literature link: \url{https://digitalcommons.unl.edu/mathfacpub/230}.
\item Additional supplied literature link: \url{https://ems.press/content/serial-article-files/46305?nt=1}.
\item Literature assessment dated exactly 2026-08-21: the supplied curation labels the problem partially solved and reports proofs for substantial configurations, including a cited line-count regime, while describing the general finite-set statement as unsettled. This is supplied dated curation, not an independently verified literature assessment. The background explicitly requests independent verification of time-sensitive status and cited claims.
\item The registry labels the statement \texttt{corrected\_verified} and reports manual semantic, typographic, and LaTeX review. These are supplied editorial assertions, not independent mathematical verification of this draft.
\end{itemize}
\end{document}
